\begin{table}[htbp]
\centering
\caption{Standardized Distributional Effect (SDE)}
\label{tab:sde}
\begin{tabular}{lcccccccc}
\hline\hline
Outcome & $\hat{\beta}$ & SE & SD($Y$) & SDE & SE(SDE) & Classification & $N$ \\
\hline
Drug arrest rate & -0.748 & 1.825 & 32.54 & -0.0230 & 0.0561 & Small & 18,192 \\
\hline\hline
\end{tabular}
\begin{tablenotes}[flushleft]
\small
\item \textbf{Country:} United States. \textbf{Research question:} Does Canada's cannabis legalization generate drug enforcement spillovers in US border counties? \textbf{Policy mechanism:} Cross-border drug policy asymmetry creates arbitrage incentives that may increase drug-market enforcement activity near border crossings. \textbf{Outcome definition:} Drug/narcotic-related arrests per 100,000 reporting population (UCR). \textbf{Treatment:} Canadian Cannabis Act effective October 17, 2018. \textbf{Data:} UCR Arrests by Age, Sex, and Race (Kaplan, Harvard Dataverse), 2014--2023, restricted to eight prohibition states. \textbf{Method:} Difference-in-differences comparing border and interior counties with three-regime interactions (post-legalization, COVID border closure, post-reopening), county and state-by-quarter FE, state-clustered SEs. \textbf{Sample:} 18,192 county-quarter observations from 457 counties (51 border) in ID, MT, ND, MN, OH, PA, NY, NH. Classification refers to magnitude, not statistical significance.
\end{tablenotes}
\end{table}
